\section{Métodos de Solução para Programação Estocástica Multiestágio}
\label{sec:rl_sddp}

A extensão natural de dois para múltiplos estágios introduz um desafio computacional central: a abordagem mais direta de solução, a Formulação Determinística Equivalente (DEQ), que representa explicitamente toda a árvore de cenários em um único programa linear, tem porte que cresce exponencialmente com o número de estágios e ramos, conforme discutido na Seção~\ref{sec:prog_estocastica}. Para horizontes de doze meses com múltiplos cenários por período, a DEQ rapidamente excede os limites de memória disponíveis, tornando-se computacionalmente inviável.

O método de Programação Dinâmica Dual Estocástica (PDDE), introduzido por \citeonline{pereira1991} para o planejamento de operação de sistemas hidrotérmicos, contorna esse problema ao substituir a função de custo futuro por uma aproximação linear por partes construída iterativamente a partir de cortes de Benders. O tamanho de cada subproblema de estágio permanece fixo e independente do horizonte total, tornando o método escalável para instâncias de grande porte.

\cite{shapiro2011} analisa formalmente as propriedades estatísticas e a convergência do PDDE, estabelecendo as condições sob as quais os cortes acumulados aproximam a função de valor ótima. \cite{shapiro2013} estendem o método para formulações com aversão a risco, incorporando medidas como CVaR na função objetivo, e apresentam estudos computacionais relacionados ao planejamento do sistema interligado brasileiro, demonstrando a aplicabilidade do método em escala real. \cite{homemdemello2011} contribuem com uma análise prática sobre estratégias de amostragem e critérios de parada para o PDDE, avaliando o compromisso entre qualidade da solução e custo computacional no contexto do despacho hidrotérmico de longo prazo.

\cite{philpott2012} incorporam medidas de risco coerentes ao PDDE e comparam políticas com diferentes níveis de aversão ao risco em um sistema hidrotérmico, mostrando que a escolha da medida de risco afeta significativamente as decisões de despacho e o perfil de exposição do agente. \cite{street2020} tratam de uma questão de modelagem relevante para a estrutura de subproblemas do PDDE: a ordem entre a tomada de decisão e a revelação da incerteza dentro de cada período. Os autores mostram que a simplificação \emph{Hazard-Decision}, em que as decisões são tomadas assumindo conhecimento perfeito das afluências do período corrente, introduz distorções econômicas mensuráveis no despacho hidrotérmico brasileiro, reforçando a importância de modelar explicitamente a precedência entre decisão e realização da incerteza.

Do ponto de vista computacional, \citeonline{dowson2021} apresentam o \texttt{SDDP.jl}, pacote em Julia que implementa o PDDE sobre o ambiente de modelagem JuMP. O pacote organiza o problema como um grafo de política, uma generalização da árvore de cenários que acomoda estruturas mais flexíveis, incluindo a decomposição de cada período em nós de decisão e nós de liquidação, e gerencia a geração de cortes, os passes \emph{forward} e \emph{backward} e o cálculo dos limites de convergência. É a ferramenta utilizada na implementação da PDDE neste trabalho.

Em síntese, a PDDE é uma técnica consolidada para problemas de energia com estrutura multiestágio, com aplicações que vão do despacho hidrotérmico ao gerenciamento de risco com medidas coerentes. No entanto, sua aplicação à decisão sequencial de compra e venda de contratos bilaterais por uma comercializadora, com comparação sistemática frente ao DEQ em diferentes portes de instância, permanece pouco explorada na literatura.
